% ============================================================================
%  D/00-map.tex — bản đồ phần D
%  Ngày tạo: 2026-09-07 | Sửa gần nhất: 2026-09-07 | Ôn gần nhất: chưa
% ============================================================================
\documentclass[../english.tex]{subfiles}
\begin{document}
\section*{Bản đồ phần D --- Sắc thái}

\begin{tomtat}
Ba chương, ba lớp sắc thái, và chúng độc lập hơn các phần khác --- đọc theo thứ tự nào cũng được. Nhưng nếu chọn một chương để đọc kỹ, hãy chọn chương 17: kết hợp từ là thứ cho hiệu quả rõ nhất trên cả đọc lẫn viết, và cũng là thứ bị bỏ sót nhiều nhất trong cách học truyền thống.
\end{tomtat}

\subsection*{A. Bảng tra ngược: bạn đang gặp gì \(\rightarrow\) đọc chương nào}
\begin{center}\footnotesize
\begin{tabularx}{\linewidth}{@{}L{5.6cm}L{2.0cm}Y@{}}
\toprule
\textbf{Tình huống} & \textbf{Chương} & \textbf{Công cụ chính}\\
\midrule
Từ điển cho ba từ, không biết chọn cái nào & 16 & Văn vực, bốn chiều sắc thái, tra dữ liệu thật\\
Câu đúng ngữ pháp mà người bản ngữ nói ``lạ'' & 17 & Kết hợp từ, mẫu câu, học theo cụm\\
Hiểu từng chữ mà không nắm được cách nghĩ của ngành & 18 & Ẩn dụ nền, khung tư duy, ẩn dụ chết\\
\bottomrule
\end{tabularx}
\end{center}

\subsection*{B. Vì sao chia như vậy}
Ba chương này ứng với ba mức mà nghĩa được quyết định ngoài phạm vi định nghĩa từ điển. Ở mức \emph{từ}, nghĩa được quyết bởi các từ khác mà người nói đã \emph{không} chọn --- đó là văn vực và sắc thái. Ở mức \emph{tổ hợp}, nghĩa và tính tự nhiên được quyết bởi thói quen kết hợp của cộng đồng, không bởi logic. Ở mức \emph{hệ khái niệm}, cách một ngành ẩn dụ hoá đối tượng của nó quyết định câu hỏi nào được coi là hợp lý.

Cách chia này khác với cách chia thông thường của sách từ vựng --- theo chủ đề (kinh tế, môi trường, giáo dục) --- và khác có chủ ý. Chia theo chủ đề cho bạn danh sách từ; chia theo \emph{cơ chế} cho bạn cách tự xử lý mọi từ mới bạn gặp về sau.

\subsection*{C. Phụ thuộc và thứ tự đọc}
\begin{figure}[H]\centering
\begin{tikzpicture}[node distance=0.58cm and 1.0cm,
  m/.style={draw=steel,fill=bluebg,rounded corners=2pt,inner sep=4.5pt,align=center,font=\small,text width=2.9cm},
  o/.style={draw=accent,fill=amberbg,rounded corners=2pt,inner sep=4.5pt,align=center,font=\small,text width=2.9cm},
  ar/.style={-{Latex[length=2.2mm]},draw=steel,thick},
  ard/.style={-{Latex[length=2.2mm]},draw=tiercol,thick,dashed},
  lb/.style={font=\scriptsize\itshape,color=reddark,align=center,text width=2.4cm}]
\node[m] (a2) {\ptr{A/01}--\ptr{A/02}\\ nghĩa nằm ở cách dùng};
\node[m,right=of a2] (c16) {\textbf{16.} Gần nghĩa và văn vực};
\node[o,right=of c16] (c17) {\textbf{17.} Kết hợp từ};
\node[m,below=1.0cm of c17] (c18) {\textbf{18.} Ẩn dụ và khung tư duy};
\draw[ar] (a2) -- (c16);
\draw[ar] (c16) -- node[lb,above=0pt] {từ \(\rightarrow\) tổ hợp} (c17);
\draw[ard] (c17) -- node[lb,right=0pt] {tổ hợp \(\rightarrow\) hệ khái niệm} (c18);
\draw[ard] (a2.south) .. controls +(0,-1.1) and +(-1.0,0) .. (c18.west);
\end{tikzpicture}
\caption{Ba mức mà nghĩa được quyết định ngoài định nghĩa từ điển: từ, tổ hợp, và hệ khái niệm. Chương 17 tô hổ phách vì nó cho hiệu quả rõ nhất trên cả đọc lẫn viết.}
\label{fig:map-D-en}
\end{figure}

\subsection*{D. Cốt lõi hay tham khảo}
\textbf{Cốt lõi:} bốn chiều sắc thái và cách tra bằng dữ liệu thật ở chương 16; toàn bộ chương 17; và mục về ẩn dụ trong chính ngành bạn ở chương 18. \textbf{Tham khảo:} phần phân loại chi tiết các kiểu kết hợp từ ở chương 17 và phần lý thuyết về ẩn dụ ý niệm ở chương 18 --- đọc để biết cơ chế, rồi quay lại khi cần.
\end{document}
